% ============================================================
% O(delta^2) Substrate-Locality Umbrella Theorem at the Host
%   Vertex of the 600-Cell Edge Graph at Vertex-Aligned Reading C
% Umbrella-Assembly Registration of THEO-DSL-5 Candidate alpha_2 =
%   -9/phi^2 at Sketch-Document Layer 3 Under Framework-Extension
%   Hypothesis (H5)
% Conscious Point Physics Flagship Paper Series
%   (F.1 Dynamical Substrate Law sub-question hardened-theorem artifact;
%    ELEVENTH and FINAL artifact in the F.1 Layer 3 hardened-theorem
%    sequence; SIXTH and FINAL artifact in OPEN-FP-F1-1 Route C sequence
%    2; closes OPEN-FP-F1-1 at sketch-document Layer 3 rigor.)
%
% Patch 0592 (Session 145, 26 May 2026) -- sixth and final OPEN-FP-F1-1
%   Route C sequence 2 closure Patch; the umbrella-assembly artifact
%   registering the THEO-DSL-5 candidate alpha_2 = -9/phi^2 from
%   Patch 0591 into the theorem-registry SD section, framing the
%   cross-sector consistency with Capotauro v2.0, and registering the
%   OPEN-FP-F1-1 closure at sketch-document Layer 3 rigor with explicit
%   open registration of the publication-grade-promotion roadmap (DG-1
%   sub-options A/B/C).
% ============================================================
% Version 1.0 --- 17 August 2026 (Patch 3212):
%   added the generated plain-language appendix glossing the internal
%   programme identifiers used in this paper (founder jargon ruling, 16
%   Aug 2026). No claim, equation, number or section changed.
%   This is the FIRST version stamp assigned to this paper; 1.0 denotes
%   the first stamped revision, NOT a claim of v1.0 shipped grade.

\documentclass[11pt,letterpaper]{article}
\usepackage[utf8]{inputenc}
\usepackage[margin=1in]{geometry}
\usepackage[T1]{fontenc}
\usepackage{amsmath,amssymb,amsfonts,amsthm}
\usepackage{xcolor}
\usepackage{hyperref}
\usepackage{enumitem}
\usepackage{parskip}
\usepackage{mdframed}

\hypersetup{
    colorlinks=true,
    linkcolor=blue,
    citecolor=blue,
    urlcolor=blue
}

\newtheorem{theorem}{Theorem}[section]
\newtheorem{proposition}[theorem]{Proposition}
\newtheorem{lemma}[theorem]{Lemma}
\newtheorem{corollary}[theorem]{Corollary}
\newtheorem{definition}[theorem]{Definition}
\newtheorem{remark}[theorem]{Remark}

\newcommand{\nhat}{\hat{n}}
\newcommand{\vhost}{v_{\text{host}}}
\newcommand{\Gsixcell}{\Gamma_{600}}
\newcommand{\Reals}{\mathbb{R}}
\newcommand{\Hthree}{H_3}
\newcommand{\Hfour}{H_4}
\newcommand{\Icoh}{I_h}
\newcommand{\jvec}{\vec{j}}

\title{$\mathcal{O}(\delta^2)$ Substrate-Locality Umbrella Theorem at the Host Vertex \\
of the 600-Cell Edge Graph at Vertex-Aligned Reading C \\
\large Umbrella-Assembly Registration of THEO-DSL-5 Candidate \\
$\alpha_2 = -9/\phi^2$ at Sketch-Document Layer 3 Under (H5) \\
(OPEN-FP-F1-1 Route C Sequence 2 Final Closure Artifact)}

\author{Thomas Lee Abshier, ND, and Claude Opus 4.7 \\ Hyperphysics Institute --- F.1 Dynamical Substrate Law trajectory}

\date{26 May 2026 (Session 145, CPP Patch 0592)\\[2pt]
{\small Version~1.0 --- 17 August 2026 (Patch 3212: added the generated plain-language appendix glossing the internal programme identifiers used in this paper (founder jargon ruling, 16 Aug 2026). No claim, equation, number or section changed.)}}

\begin{document}
\maketitle

\begin{abstract}
We assemble the eleven F.1 Dynamical Substrate Law Layer~3 hardened-theorem artifacts (Patches~0550 + 0551 + 0552 + 0571 + 0585 + 0587 + 0588 + 0589 + 0590 + 0591 plus the present umbrella Patch~0592) into a single umbrella theorem registering the THEO-DSL-5 candidate $\alpha_2 = -9/\phi^2 = -9(2 - \phi) \approx -3.438$ at sketch-document Layer~3 under framework-extension hypothesis~(H5) per Session~144 close handover decision-gate bundle DG-1 sub-option~(A). The umbrella statement asserts: under hypotheses (H1)~vertex-aligned Reading~C with $\nhat = \vhost$, (H2)~600-cell unit-vertex normalisation, (H3)~icosahedral residual symmetry $\Hthree = \Icoh$, (H4)~Mechanism~A's rate function $r(\hat{e}) = r_0(1 + \delta\,\hat{e} \cdot \nhat)$, and (H5)~the path-class weight ansatz $W(P) = 1$ for $P \in \mathcal{P}_2(\vhost)$ at sketch-document Layer~3, the second-order substrate-current $\jvec_2(\vhost) = \alpha_2 \nhat \cdot r_0 \delta^2$ has scalar coefficient $\alpha_2 = -9/\phi^2$ in closed form. The umbrella statement closes \textbf{OPEN-FP-F1-1} at sketch-document Layer~3 rigor; publication-grade Layer~3 hardening of THEO-DSL-5 requires DG-1 sub-option~(A)~path-integral rigorisation, (B)~new framework axiom~MA.3, or (C)~Layer~4 derivation~OPEN-FP-F1-2. Cross-sector consistency with Capotauro~v2.0~\S5--\S6 is structurally parallel under matched vertex-aligned Reading~C conventions: the closed-form sign + golden-ratio-magnitude structure of $\alpha_2 = -9/\phi^2$ in F.1's temporal sector mirrors the analogous second-order substrate-locality correction in Capotauro's spatial sector. With Patch~0592 in hand, the F.1 Layer~3 hardened-theorem sequence stands at \textbf{eleven artifacts} and the Route~C sequence~2 trajectory is COMPLETE within the originally-estimated 4--8 session budget (six sessions total: Patches~0587 + 0588 + 0589 + 0590 + 0591 + 0592). The forward queue post-Patch-0592 returns to the post-v1.0 priority queue established at F.1~v1.0~SHIP (Session~54): (A)~OPEN-FP-SF-4-1 Picture~A, (B)~SM-5 cooperation, (C)~SF-2 for $\delta_{CP}$, (D)~anthology, (E)~TATWD book project, (F)~JUNO peer-review update, with OPEN-FP-F1-1-publication-grade-hardening (via DG-1 sub-option~A/B/C) added to the priority queue at a future cycle.
\end{abstract}

\section{Introduction and the role of the umbrella artifact}\label{sec:introduction}

This paper assembles the eleven F.1 Dynamical Substrate Law Layer~3 hardened-theorem artifacts into a single umbrella theorem registering the THEO-DSL-5 candidate at sketch-document Layer~3 under framework-extension hypothesis~(H5). The umbrella's role is registration and scope-framing, not substantive new derivation: all substantive content was established in Patches~0550--0591 across the prior eleven sessions of the Route~C trajectory. Patch~0592 closes the OPEN-FP-F1-1 trajectory by:
\begin{enumerate}[label=(\roman*),leftmargin=2em]
\item Stating the umbrella THEO-DSL-5 candidate theorem at sketch-document Layer~3 rigor, citing the chain of supporting artifacts.
\item Registering the rigor distinction between the publication-grade Layer~3 unconditional supporting artifacts (the geometric-primitive + structural-theorem set: G1 + G1.1 + G1.2 + G2 + G2.1 + G2.3 + G2.4 + structural theorem at Patches~0551 + 0552 + 0571 + 0587 + 0588 + 0589 + 0590 + 0585) and the sketch-document Layer~3 path-class enumeration under (H5) at Patch~0591.
\item Framing the cross-sector consistency with Capotauro~v2.0 at structural parallelism level.
\item Registering the publication-grade-promotion roadmap (DG-1 sub-options~A/B/C) for downstream consumers.
\item Closing OPEN-FP-F1-1 at sketch-document Layer~3 rigor and returning the forward queue to the F.1~v1.0~SHIP post-v1.0 priority structure.
\end{enumerate}

\paragraph{Position in the F.1 Layer 3 hardened-theorem sequence (now eleven artifacts).}
The F.1 Dynamical Substrate Law substrate-locality programme has produced ten publication-grade Layer~3 hardened-theorem artifacts plus one sketch-document Layer~3 path-class enumeration artifact prior to the present umbrella:
\begin{itemize}[leftmargin=2em]
\item \textbf{$\mathcal{O}(\delta^1)$ first-shell building-block trio} (publication-grade Layer~3 unconditional): Patches~0550 (perturbation-locality propagation) + 0551 (G1.2 first-shell perpendicularity) + 0552 (G1.1 host-to-first-shell projection) + 0571 (G1 first-shell inner-product primitive).
\item \textbf{All-orders structural theorem} (publication-grade Layer~3 unconditional): Patch~0585 (second-order parallel-to-$\nhat$ all-orders structural theorem).
\item \textbf{Second-shell building-block tetrad} (publication-grade Layer~3 unconditional): Patches~0587 (G2 second-shell inner-product primitive) + 0588 (G2.1 host-to-second-shell projection) + 0589 (G2.4 second-shell perpendicularity) + 0590 (G2.3 first-shell-to-second-shell projection).
\item \textbf{Path-class enumeration} (sketch-document Layer~3 under (H5)): Patch~0591 ($\mathcal{O}(\delta^2)$ path-class weight enumeration yielding $\alpha_2 = -9/\phi^2$).
\item \textbf{Umbrella registration} (sketch-document Layer~3, present Patch~0592): the THEO-DSL-5 candidate registration.
\end{itemize}
With Patch~0592, the F.1 Layer~3 hardened-theorem sequence has \textbf{eleven artifacts}.

\paragraph{Sequence-2 budget closure.}
The Route~C sequence~2 trajectory has consumed \textbf{six sessions} (Patches~0587 + 0588 + 0589 + 0590 + 0591 + 0592) of the original 4--8 session budget. The budget closure is at the central value (6 of 4--8 estimated), confirming the original scoping estimate. The Route~C sequence~1 trajectory (Patch~0585's all-orders structural theorem) consumed one session prior, giving a Route~C total of seven sessions for the OPEN-FP-F1-1 closure at sketch-document Layer~3 rigor.

\section{The umbrella theorem}\label{sec:umbrella_theorem}

\subsection{Hypotheses}\label{sec:umbrella_hypotheses}

The umbrella theorem operates under five hypotheses, the first four of which are publication-grade-rigorisable and the fifth of which is at sketch-document Layer~3 rigor:

\begin{itemize}[leftmargin=2em]
\item \textbf{(H1)} Vertex-aligned Reading~C: $\nhat = \vhost$, where $\vhost \in V(\Gsixcell)$ is the host vertex and $\nhat$ is the substrate-locality reading-axis. (Publication-grade Layer~3 hardening in F.1~v1.0~\S3 + Patch~0550 \S2.)
\item \textbf{(H2)} 600-cell unit-vertex normalisation: every $v \in V(\Gsixcell)$ has $|v| = 1$. (Publication-grade Layer~3 hardening; standard polytope-theoretic normalisation.)
\item \textbf{(H3)} Icosahedral residual symmetry: the residual point-group at $\vhost$ in the 600-cell symmetry group $\Hfour$ is $\Hthree = \Icoh$ of order 120. (Publication-grade Layer~3 hardening in Patch~0571 + standard polytope theory; \cite{coxeter1973}.)
\item \textbf{(H4)} Mechanism~A's rate function: $r(\hat{e}) = r_0(1 + \delta\,\hat{e} \cdot \nhat)$ on every directed edge $\hat{e} \in \Reals^4$, $|\hat{e}| = 1$. (Publication-grade Layer~3 hardening in F.1~v1.0 \S4 axiom MA.1.)
\item \textbf{(H5)} \emph{Framework-extension at sketch-document Layer~3}: the $\mathcal{O}(\delta^2)$ substrate-current at $\vhost$ takes the path-class sum form $\jvec_2(\vhost) = r_0 \delta^2 \sum_{P \in \mathcal{P}_2(\vhost)} W(P)\,\omega(P)\,\hat{e}_1(P)$ with $W(P) = 1$ uniformly across $P \in \mathcal{P}_2(\vhost)$, where $\omega(P) = (\hat{e}_1 \cdot \nhat)(\hat{e}_2 \cdot \nhat)$ is the path-class direction-product weight and $\hat{e}_1(P), \hat{e}_2(P)$ are the first- and second-edge unit-direction vectors of path $P$. (Sketch-document Layer~3 per DG-1 sub-option~(A); publication-grade hardening requires DG-1 sub-options~A/B/C resolution.)
\end{itemize}

\subsection{Statement of THEO-DSL-5 candidate}\label{sec:theorem_statement}

\begin{theorem}[THEO-DSL-5 umbrella candidate; $\mathcal{O}(\delta^2)$ substrate-current at $\vhost$]\label{thm:theo_dsl_5}
Under hypotheses (H1)--(H5), the second-order substrate-current at the host vertex of the 600-cell edge graph at vertex-aligned Reading~C takes the closed-form value
\begin{equation}\label{eq:theo_dsl_5}
\jvec_2(\vhost) \;=\; \alpha_2 \,\nhat \cdot r_0 \delta^2, \qquad \alpha_2 \;=\; -\frac{9}{\phi^2} \;=\; -9(2 - \phi) \;\approx\; -3.438.
\end{equation}
The parallel-to-$\nhat$ direction follows from Patch~0585 (all-orders structural theorem; publication-grade Layer~3 unconditional). The scalar coefficient value $-9/\phi^2$ follows from Patch~0591 (path-class weight enumeration; sketch-document Layer~3 under (H5)).
\end{theorem}

\begin{proof}[Proof of Theorem~\ref{thm:theo_dsl_5}]
The proof is a citation chain assembling the eleven F.1 Layer~3 hardened-theorem artifacts.

\textbf{Step 1: parallel-to-$\nhat$ direction.}
By Patch~0585 Theorem~1.1 (publication-grade Layer~3 unconditional under (H1)--(H3)), $\jvec_k(\vhost) = \beta_k \nhat$ for some scalar $\beta_k$ at every perturbative order $k \geq 1$; specialising to $k = 2$ gives $\jvec_2(\vhost) = \alpha_2 \nhat \cdot r_0 \delta^2$ for some scalar $\alpha_2$ to be determined.

\textbf{Step 2: scalar coefficient value.}
By Patch~0591 Theorem~1.1 (sketch-document Layer~3 under (H1)--(H5)), the scalar $\alpha_2$ takes the closed-form value $\alpha_2 = -9/\phi^2$ under the path-class weight ansatz $W(P) = 1$ specified in (H5).

\textbf{Step 3: chain integrity.}
The path-class enumeration of Patch~0591 imports seven geometric-primitive results plus one inline lemma at publication-grade Layer~3 unconditional rigor: G1 (Patch~0571), G1.1 (Patch~0552), G1.2 (Patch~0551), G2 (Patch~0587), G2.1 (Patch~0588), G2.3 (Patch~0590), G2.4 (Patch~0589), and Lemma~2.2 (Patch~0591's inline G3 first-shell-to-third-shell projection $-\phi/2$, at sketch-document Layer~3 via standard polytope-theoretic shell decomposition). The path-class sum's structural form is at sketch-document Layer~3 under (H5)'s $W(P) = 1$ ansatz; the geometric inputs are at publication-grade Layer~3 unconditional rigor.

Combining Steps 1--3 yields the closed-form $\alpha_2 = -9/\phi^2$ at the rigor level of the WEAKEST link in the chain, which is (H5) at sketch-document Layer~3. $\square$
\end{proof}

\subsection{Path-class decomposition: per-class contributions}\label{sec:path_class_decomposition}

For downstream consumers of THEO-DSL-5, we register the per-path-class decomposition of $\alpha_2 = -9/\phi^2$:

\begin{corollary}[Per-path-class $\nhat$-component contributions to $\alpha_2$]\label{cor:per_class_decomposition}
Under (H1)--(H5), the closed-form $\alpha_2 = -9/\phi^2$ decomposes into four path-class contributions per Patch~0591 Corollary~4.1:
\begin{itemize}[leftmargin=2em]
\item \textbf{B.1} (12 round-trip paths returning to $\vhost$): $+3/(2\phi^3) \approx +0.354$, contributing approximately $-10\%$ of $\alpha_2$ via partial cancellation.
\item \textbf{B.2} (60 first-shell-trans paths through icosahedron edges): $0$ exactly, via Patch~0551 G1.2 first-shell perpendicularity (structural vanishing).
\item \textbf{B.3} (60 cross-shell paths $\vhost \to S_1 \to S_2$): $-15/(2\phi^2) \approx -2.865$, contributing approximately $83\%$ of $\alpha_2$ (dominant).
\item \textbf{B.4} (12 third-shell paths $\vhost \to S_1 \to S_3$): $-3/(2\phi) \approx -0.927$, contributing approximately $27\%$ of $\alpha_2$ (sub-dominant, reinforcing B.3).
\end{itemize}
Total: $+3/(2\phi^3) + 0 + (-15/(2\phi^2)) + (-3/(2\phi)) = -9/\phi^2$ via golden-ratio algebra.
\end{corollary}

\begin{remark}[Structural origin of the closed-form]\label{rem:structural_origin}
The clean closed-form $-9/\phi^2$ has a transparent structural origin: the per-first-shell-vertex projection sum $S(u_i) = -3$ exactly (Patch~0591 Proposition~4.2) combined with the icosahedral sum $\sum_{i=1}^{12} \hat{u}_i = -(6/\phi)\nhat$ (an $\Icoh$-symmetry consequence) and the host-to-first-shell projection $\hat{u}_i \cdot \nhat = -1/(2\phi)$ (Patch~0552 G1.1) gives $\alpha_2 = (-1/(2\phi)) \cdot (-3) \cdot (-6/\phi) = -9/\phi^2$ in a three-factor product. The integer factor $-3$ in $S(u_i)$ emerges from golden-ratio algebraic cancellations across the four edge-class second-edge projections; the integer factor $-6$ in the icosahedral sum emerges from twice the host-to-first-shell projection times the vertex count $12 \cdot \hat{u}_i \cdot \nhat$. Both integer factors are uniform-across-the-icosahedron consequences of the $\Icoh$ symmetry at vertex-aligned Reading~C.
\end{remark}

\section{Cross-sector consistency framing with Capotauro v2.0}\label{sec:cross_sector_consistency}

\subsection{Structural parallelism at matched vertex-aligned Reading C}\label{sec:structural_parallel}

Capotauro~v2.0~\S5--\S6 establishes a parallel substrate-locality structure in the spatial sector under matched vertex-aligned Reading~C conventions. The relevant cross-sector correspondence:

\begin{itemize}[leftmargin=2em]
\item \textbf{First-order coefficient correspondence.} F.1~v1.0~Theorem~7.1's $\alpha_1 = 6/\phi^2$ has a Capotauro~v2.0~\S5-derived parallel coefficient at structural-parallelism level (same golden-ratio magnitude, same sign convention under matched orientation choices, possibly different integer prefactor reflecting the temporal-vs-spatial-sector chirality-content distribution).
\item \textbf{Second-order coefficient correspondence.} The present Patch~0592's THEO-DSL-5 candidate $\alpha_2 = -9/\phi^2$ has an analogous Capotauro~v2.0~\S6-derived second-order coefficient, with parallel structural features: (a)~closed-form involving $\phi$ and small integers; (b)~sign opposite to the first-order coefficient; (c)~magnitude of order $1/\phi^2$; (d)~per-path-class decomposition mirroring the F.1 B.1--B.4 partition (with the cross-shell class dominating).
\end{itemize}

\begin{remark}[Cross-sector consistency at structural-parallelism level vs. numerical-coefficient level]\label{rem:cross_sector_levels}
The cross-sector consistency between F.1 and Capotauro~v2.0 is established at \emph{structural-parallelism level} (same closed-form classes, same sign patterns, same golden-ratio order of magnitude) but NOT at strict numerical-coefficient level (the integer prefactors $-9$ in F.1's $\alpha_2$ and the analogous Capotauro coefficient may differ depending on the precise chirality-content normalisation in the two sectors). A future cross-sector consistency Patch executing the parallel Capotauro~v2.0 path-class enumeration at $\mathcal{O}(\delta^2)$ would establish strict numerical-coefficient cross-sector matching or register a discrepancy as a substantive cross-sector consistency target. The structural parallelism is sufficient for the present umbrella registration.
\end{remark}

\subsection{Capotauro v2.0 W-bracelet protection parallel}\label{sec:w_bracelet_parallel}

The Capotauro~v2.0 W-bracelet protection structure (Capotauro v2.0~\S5--\S6, chirality-content protection between first and second shells in the spatial sector) corresponds to the F.1 second-shell-internal perpendicularity at Patch~0589 (G2.4 second-shell-to-second-shell edge-direction perpendicularity yielding rate-asymmetry vanishing). Both reflect the same underlying mechanism: \emph{within-shell perpendicularity at every shell yields zero rate-asymmetry contribution at the relevant perturbative order, providing a structural protection against contamination of higher-shell coefficients by lower-shell artifacts}. The Patch~0591 Corollary~4.2 shell-perpendicularity propagation pattern documents this in the F.1 sector; the Capotauro~v2.0 W-bracelet protection registers the same in the chirality-content sector.

\section{OPEN-FP-F1-1 closure and the publication-grade-promotion roadmap}\label{sec:open_fp_f1_1_closure}

\subsection{OPEN-FP-F1-1 closure status}\label{sec:closure_status}

\begin{remark}[OPEN-FP-F1-1 closure at sketch-document Layer 3 rigor]\label{rem:open_fp_f1_1_closure}
With Patch~0592 in hand, \textbf{OPEN-FP-F1-1 is CLOSED at sketch-document Layer 3 rigor under (H1)--(H5)}. The structural sub-question (parallel-to-$\nhat$ direction at all orders) was closed at publication-grade Layer~3 unconditional rigor via Patch~0585. The coefficient sub-question (closed-form value $\alpha_2 = -9/\phi^2$) is closed at sketch-document Layer~3 rigor under (H5). The combined status is sketch-document Layer~3 closure of OPEN-FP-F1-1 at the rigor of the weakest link, i.e., (H5).
\end{remark}

\subsection{Publication-grade-promotion roadmap}\label{sec:promotion_roadmap}

For downstream consumers requiring publication-grade Layer~3 unconditional hardening of THEO-DSL-5, three sub-option routes remain (per DG-1 sub-options~A/B/C of the Session~144 close handover):

\begin{itemize}[leftmargin=2em]
\item \textbf{DG-1 sub-option (A) --- Path-integral rigorisation.} Reformulate (H5)'s path-class weight ansatz $W(P) = 1$ as a derived consequence of a publication-grade path-integral construction over the 600-cell substrate dynamics. The path-class enumeration of Patch~0591 would then become a corollary of the path-integral derivation rather than a postulate at sketch-document Layer~3. This is the LIGHTEST-WEIGHT promotion route: it preserves the present Patches~0585--0591 chain unchanged and adds a rigorisation layer above (H5).
\item \textbf{DG-1 sub-option (B) --- New framework axiom MA.3.} Introduce a new Mechanism~A axiom MA.3 explicitly registering the path-class weight $W(P) = 1$ as a framework primitive. Publication-grade Layer~3 unconditional rigor of THEO-DSL-5 follows immediately. The trade-off: a new axiom (loss of axiomatic parsimony) in exchange for promotion.
\item \textbf{DG-1 sub-option (C) --- Layer 4 derivation OPEN-FP-F1-2.} Derive Mechanism~A's second-order path-class structure (including $W(P) = 1$) from CPP primitive axioms~$A1$--$A11$ via a substrate-derivation programme paralleling the F.1 v1.0 Layer~4 derivation of the substrate-locality framework. This is the HEAVIEST-WEIGHT promotion route but yields the strongest result: publication-grade Layer~4 derivation of THEO-DSL-5 from CPP primitives.
\end{itemize}

Each route is downstream of the present Route~C sequence~2 trajectory and would be a multi-Patch endeavour at a future cycle. Selection among the three routes is a strategic decision for downstream prioritisation; (A) is recommended as the default unless one of (B) or (C) becomes attractive on independent grounds (e.g., (C) for foundational unification at Layer~4 rigor, or (B) for minimal-effort registration of the framework primitive).

\subsection{Anti-erasure registration: load-bearing nature of (H5)}\label{sec:anti_erasure}

\begin{remark}[Sketch-document Layer 3 rigor is load-bearing]\label{rem:H5_load_bearing}
The (H5) framework-extension hypothesis at sketch-document Layer~3 is the LOAD-BEARING scope qualifier on Theorem~\ref{thm:theo_dsl_5}'s rigor level. Downstream consumers of $\alpha_2 = -9/\phi^2$ should treat the value as a \emph{candidate result pending publication-grade-promotion via DG-1 sub-option~A/B/C resolution}, NOT as a publication-grade Layer~3 unconditional prediction. Reports, papers, or external communications citing THEO-DSL-5 should include the qualifier ``sketch-document Layer~3 under (H5)'' or equivalent. The publication-grade-promotion roadmap of \S\ref{sec:promotion_roadmap} provides three clear routes for downstream consumers who require publication-grade rigor; the present umbrella does not pre-select among these routes.
\end{remark}

\section{Forward queue and return to post-v1.0 priority structure}\label{sec:forward_queue}

\subsection{Closure of the Route C trajectory}\label{sec:route_c_closure}

With Patch~0592, the Route~C trajectory of OPEN-FP-F1-1 closure is COMPLETE:
\begin{itemize}[leftmargin=2em]
\item Route~C sequence~1: closed at Patch~0585 (1 session).
\item Route~C sequence~2: closed at Patches~0587--0592 (6 sessions).
\item Route~C total: 7 sessions.
\end{itemize}
The Route~C closure is within the original 4--8 session budget for sequence~2 (at the central value 6 of 4--8) and matches the F.1 v1.0~SHIP-cycle scoping estimate.

\subsection{Return to F.1 v1.0 SHIP post-v1.0 priority queue}\label{sec:post_v1_priority_queue}

With OPEN-FP-F1-1 closed at sketch-document Layer~3 rigor, the forward queue returns to the post-v1.0 priority queue established at F.1 v1.0~SHIP (Session~54, 23~May~2026):
\begin{itemize}[leftmargin=2em]
\item \textbf{(A) OPEN-FP-SF-4-1 Picture A}: neutrino sector flagship completion; primary post-v1.0 priority.
\item \textbf{(B) SM-5 cooperation}: ChatGPT-led SM-5 reviewer cooperation.
\item \textbf{(C) SF-2 for $\delta_{CP}$}: CP-violation phase in the neutrino sector.
\item \textbf{(D) Anthology}: omnibus paper-series anthology compilation.
\item \textbf{(E) TATWD book project}: ``Tetrahedrons All the Way Down'' book project.
\item \textbf{(F) JUNO peer-review update}: JUNO experimental-results update for SF-4.
\item \textbf{(G) NEW: OPEN-FP-F1-1 publication-grade hardening} via DG-1 sub-option~A/B/C resolution; added to the priority queue at the present cycle close. Default route: sub-option~(A) path-integral rigorisation.
\end{itemize}

Priority selection between (A)--(G) is a strategic decision for the next session-open handover.

\section{Scope, framework-locality, and explicit exclusions}\label{sec:scope_and_exclusions}

The umbrella theorem inherits the scope and exclusions of its constituent artifacts. Explicitly enumerating:

\begin{enumerate}[label=\textbf{(E\arabic*)}]
\item \textbf{(H5) framework-extension at sketch-document Layer 3.} See \S\ref{sec:promotion_roadmap} for the publication-grade-promotion roadmap. The closed-form $\alpha_2 = -9/\phi^2$ is rigor-bounded by (H5).

\item \textbf{Higher-order extensions $\mathcal{O}(\delta^k)$ for $k \geq 3$ are NOT addressed.} Path-class enumeration at $\mathcal{O}(\delta^3)$ would require additional geometric primitives (G3 + G3.0 + G3.1 + etc.) plus extension of (H5)'s path-class weight ansatz to 3-edge paths. Registered as a candidate future OPEN-FP-F1-N target if higher-order extensions become a programme priority.

\item \textbf{Non-vertex-aligned Reading C variants are NOT covered.} The theorem requires $\nhat = \vhost$ (hypothesis (H1)). Registered as OPEN-FP-F1-5.

\item \textbf{Layer 4 derivation of the 600-cell substrate from CPP primitive axioms.} OPEN-FP-F1-2; one of the three promotion routes per \S\ref{sec:promotion_roadmap} (DG-1 sub-option~C).

\item \textbf{Cross-sector consistency with Capotauro v2.0 at strict numerical-coefficient level is NOT established.} The cross-sector consistency at structural-parallelism level is registered in \S\ref{sec:cross_sector_consistency}; strict numerical-coefficient cross-sector matching is a candidate future cross-sector consistency Patch.

\item \textbf{Mechanism A's framework-locality at $\mathcal{O}(\delta^k)$ extensions.} The Patch~0550 perturbation-locality theorem covers shell-locality at $\mathcal{O}(\delta^1)$ and is extended to $\mathcal{O}(\delta^2)$ at the present Patch level under (H5). Higher-order shell-locality extensions parallel to Patch~0550 at $\mathcal{O}(\delta^k)$ are NOT addressed.
\end{enumerate}

\section{Conclusion --- OPEN-FP-F1-1 is closed at sketch-document Layer 3}\label{sec:conclusion}

Theorem~\ref{thm:theo_dsl_5} (THEO-DSL-5 umbrella candidate) closes \textbf{OPEN-FP-F1-1} at \textbf{sketch-document Layer 3 rigor} under hypotheses (H1)--(H5). The closed-form coefficient
\begin{equation*}
\alpha_2 \;=\; -\frac{9}{\phi^2} \;=\; -9(2 - \phi) \;\approx\; -3.438
\end{equation*}
is the THEO-DSL-5 candidate value, derived via the path-class enumeration at $\mathcal{O}(\delta^2)$ assembling seven publication-grade Layer~3 unconditional geometric primitives plus the all-orders structural theorem under the (H5) path-class weight ansatz $W(P) = 1$. Publication-grade Layer~3 unconditional hardening of THEO-DSL-5 requires DG-1 sub-option~A/B/C resolution per \S\ref{sec:promotion_roadmap}.

\paragraph{F.1 Layer 3 hardened-theorem sequence at eleven artifacts.}
With Patch~0592, the F.1 Dynamical Substrate Law Layer~3 hardened-theorem sequence is at its target completeness for the Route~C trajectory: eleven artifacts spanning Patches~0550 + 0551 + 0552 + 0571 + 0585 + 0587 + 0588 + 0589 + 0590 + 0591 + 0592.

\paragraph{Forward queue.}
The forward queue returns to the F.1 v1.0~SHIP post-v1.0 priority structure with the new addition of OPEN-FP-F1-1 publication-grade hardening (via DG-1 sub-option~A/B/C). Priority selection among (A) OPEN-FP-SF-4-1 Picture~A, (B) SM-5 cooperation, (C) SF-2 for $\delta_{CP}$, (D) anthology, (E) TATWD book project, (F) JUNO peer-review update, and (G) OPEN-FP-F1-1 publication-grade hardening is a strategic decision for the next session-open handover.

\paragraph{Anti-erasure note.}
The present umbrella preserves: (a)~the explicit registration of (H5) as load-bearing sketch-document Layer~3 (Remark~\ref{rem:H5_load_bearing}); (b)~the three-route publication-grade-promotion roadmap (\S\ref{sec:promotion_roadmap}) with default recommendation but no pre-selection among A/B/C; (c)~the cross-sector consistency framing at structural-parallelism level explicitly distinguished from strict numerical-coefficient level (Remark~\ref{rem:cross_sector_levels}); (d)~the per-path-class decomposition (Corollary~\ref{cor:per_class_decomposition}) registering the dominant cross-shell contribution and the partial-cancellation round-trip contribution, so that downstream consumers understand the structural composition rather than treating $-9/\phi^2$ as a single opaque value; (e)~the chain-rigor-level analysis in Theorem~\ref{thm:theo_dsl_5}'s proof Step~3 (weakest-link rigor analysis explicit rather than buried).

\bibliographystyle{plain}
\begin{thebibliography}{99}

\bibitem{coxeter1973}
H.~S.~M.~Coxeter,
\textit{Regular Polytopes},
3rd ed., Dover Publications, New York, 1973.

\bibitem{humphreys1990}
J.~E.~Humphreys,
\textit{Reflection Groups and Coxeter Groups},
Cambridge Studies in Advanced Mathematics 29, Cambridge University Press, 1990.

\bibitem{abshier_f1_dynamical_substrate_law}
T.~L.~Abshier,
``The Dynamical Substrate Law: Substrate-Locality of DI-Bit Currents at Vertex-Aligned Reading~C in the 600-Cell,''
CPP Flagship Paper Series (F-line), Version 1.0, 24~May~2026, OSF DOI \texttt{10.17605/OSF.IO/JXE8D}.

\bibitem{abshier_capotauro_v2}
T.~L.~Abshier,
``The Capotauro Framework v2.0: Spatial-Sector Substrate-Locality and W-Bracelet Chirality-Content Protection,''
CPP Flagship Paper Series, prior to F.1 v1.0~SHIP, 22~May~2026.

\bibitem{abshier_perturbation_locality}
T.~L.~Abshier, ``Perturbation-Theory Propagation Rule and Shell-Locality,''
CPP F.1 hardened-theorem artifact, Patch~0550.

\bibitem{abshier_first_shell_perpendicularity}
T.~L.~Abshier, ``First-Shell-to-First-Shell Edge-Direction Perpendicularity (G1.2),''
CPP F.1 hardened-theorem artifact, Patch~0551.

\bibitem{abshier_first_shell_projection}
T.~L.~Abshier, ``Host-to-First-Shell Uniform Projection (G1.1),''
CPP F.1 hardened-theorem artifact, Patch~0552.

\bibitem{abshier_g1_hardening}
T.~L.~Abshier and Claude Opus 4.7, ``First-Shell Inner-Product Primitive (G1),''
CPP F.1 hardened-theorem artifact, Patch~0571.

\bibitem{abshier_structural_theorem}
T.~L.~Abshier and Claude Opus 4.7, ``Second-Order Parallel-to-$\hat{n}$ Structural Theorem,''
CPP F.1 hardened-theorem artifact, Patch~0585.

\bibitem{abshier_g2_hardening}
T.~L.~Abshier and Claude Opus 4.7, ``Second-Shell Inner-Product Primitive (G2),''
CPP F.1 hardened-theorem artifact, Patch~0587.

\bibitem{abshier_g21_hardening}
T.~L.~Abshier and Claude Opus 4.7, ``Host-to-Second-Shell Uniform Projection (G2.1),''
CPP F.1 hardened-theorem artifact, Patch~0588.

\bibitem{abshier_g24_hardening}
T.~L.~Abshier and Claude Opus 4.7, ``Second-Shell-to-Second-Shell Edge-Direction Perpendicularity (G2.4),''
CPP F.1 hardened-theorem artifact, Patch~0589.

\bibitem{abshier_g23_hardening}
T.~L.~Abshier and Claude Opus 4.7, ``First-Shell-to-Second-Shell Edge-Direction Projection (G2.3),''
CPP F.1 hardened-theorem artifact, Patch~0590.

\bibitem{abshier_path_class_enumeration}
T.~L.~Abshier and Claude Opus 4.7, ``Path-Class Weight Enumeration at $\mathcal{O}(\delta^2)$,''
CPP F.1 hardened-theorem artifact, Patch~0591.

\end{thebibliography}

% BEGIN GENERATED IDENTIFIER APPENDIX -- do not edit by hand
\clearpage
\section*{Appendix: Programme identifiers used in this paper}
\addcontentsline{toc}{section}{Appendix: Programme identifiers}

\noindent This programme tracks its unresolved questions under short identifiers so that a claim and its open dependencies can be cited precisely. The identifiers appearing in this paper are glossed below; no external document is needed to read them.

\begin{description}
  \item[\texttt{OPEN-FP-F1-1}] \hfill \\ Extend the F.1 substrate-locality umbrella theorem (Theorem 7.1) to \$\textbackslash\{\}mathcal\{O\}(\textbackslash\{\}delta\textasciicircum{}2)\$ by deriving the second-shell inner-product and edge-projection identities for the 600-cell, and determine whether the parallel-to-\$\textbackslash\{\}hat\{n\}\$ structure of \$\textbackslash\{\}vec\{J\}\_\{\textbackslash\{\}text\{DI-net\}\}(v\_\{\textbackslash\{\}text\{host\}\})\$ survives at second order. **Resolution**: both sub-questions closed — structural sub-question at publication-grade Layer 3 unconditional (\$\textbackslash\{\}vec\{j\}\_k(v\_\{\textbackslash\{\}text\{host\}\}) \textbackslash\{\}parallel \textbackslash\{\}hat\{n\}\$ at every \$k \textbackslash\{\}geq 1\$ via THEO-DSL-4) + coefficient sub-question at sketch-document Layer 3 under (H5) (\$\textbackslash\{\}alpha\_2 = -9/\textbackslash\{\}phi\textasciicircum{}2 = -9(2-\textbackslash\{\}phi) \textbackslash\{\}approx -3.438\$ via THEO-DSL-5 candidate; ratio \$\textbackslash\{\}alpha\_2/\textbackslash\{\}alpha\_1 = -3/2\$ exactly).
  \item[\texttt{OPEN-FP-F1-2}] \hfill \\ Derive Mechanism A (F.1 framework axioms MA.1 + MA.2: propagation-rate asymmetry \$r(\textbackslash\{\}hat\{e\}) = r\_0(1 + \textbackslash\{\}delta\textbackslash\{\},\textbackslash\{\}hat\{e\}\textbackslash\{\}cdot\textbackslash\{\}hat\{n\})\$ and framework-local current construction at \$\textbackslash\{\}mathcal\{O\}(\textbackslash\{\}delta\textasciicircum{}1)\$) from CPP primitive axioms A1–A11 alone.
  \item[\texttt{OPEN-FP-F1-5}] \hfill \\ Extend or reformulate F.1 Theorems 5.1 (host-first-shell projection), 5.2 (first-shell perpendicularity), and 7.1 (substrate-locality umbrella) under edge-aligned Reading C (\$\textbackslash\{\}hat\{n\}\$ along a 600-cell short-edge direction; residual \$D\_5\$ symmetry at edge midpoint per Patch 0596 Remark 5.1 anti-erasure correction — see below) and face-aligned Reading C (\$\textbackslash\{\}hat\{n\}\$ perpendicular to a triangular face; residual \$C\_s\$ symmetry under vertex-host convention per Patch 0597 Remark 7.1 anti-erasure correction — see below).
  \item[\texttt{OPEN-FP-F1-N}] \hfill \\ A reserved identifier for a future higher-order extension of the F.1 geometric primitives, should that become a programme priority.
  \item[\texttt{OPEN-FP-SF-4-1}] \hfill \\ Derive from CPP primitives the full mass formula \$m\_\{\textbackslash\{\}nu\_i\} = M\_0 \textbackslash\{\}cdot V\_\{\textbackslash\{\}nu\_i\}\textasciicircum{}2 \textbackslash\{\}cdot \textbackslash\{\}sigma\_\textbackslash\{\}nu\$ for unbound neutrino modes. **CLOSURE ACHIEVED v3.0**: Picture A delivers \$\textbackslash\{\}sigma\_\textbackslash\{\}nu = 1/z\textasciicircum{}\{10\}\$; α-exponent closure delivers \$V\textasciicircum{}\{1/3\} \textbackslash\{\}to 1\$ at the bound/unbound boundary; together, the full mass formula is rigorously derived from CPP axioms plus six foundational inputs (three shared with Picture A; one shared with α-exponent; three operational definitions).
\end{description}
% END GENERATED IDENTIFIER APPENDIX

\end{document}
